\section{The free chamber inverse}\label{sec:propagators}

The chamber estimates of \cite[Section~9]{QED3Manuscript} give bounds which
are uniform in the mesh, the volume, and the running mass. There are two
sources of coercivity. The internal chamber graph controls the part of a
spinor whose mean over a cluster vanishes. The averaging penalty controls
the remaining four components. The first of these bounds also identifies
the singular part of the free propagator. This permits a local regularity
estimate with constants independent of the mesh, which is needed when
the inverses on small tori are combined into an inverse on the full torus.

\subsection{The cluster decomposition}\label{subsec:cluster-decomposition}

We first allow an arbitrary mesh $a>0$, using the chamber stencil of
\cref{sec:introduction} with $\eta$ replaced by $a$. Put
$\widehat d_i(c,g)=-gn_i$ and
\[
 M_i(c,g)=\frac{w_i}{2}(\rho-\gamma\cdot\widehat d_i(c,g)).
\]

The discussion applies on the covering lattice and on a torus of side
$R\in a\N$, with any unitary character
$f(c+Re_\mu,g)=-\exp(i\theta_\mu)f(c,g)$. Pairings on the torus are
formed over a period cell. The notation $d(c,c')$ denotes the sup
distance modulo $R\Z^3$, or the covering sup distance when $R=\infty$.
They use the measure
\[
 \int\dd x\,v(x)=\frac{a^3}{48}\sum_c\sum_gv(c,g),
 \qquad
 \|f\|_2^2=\int\dd x\,|f(x)|_2^2.
\]
On four-component spinors depending only on $c$, the measure is
$a^3\sum_c$. With the cluster maps $J,J^*$ of the introduction,
set $P=JJ^*$ and $Q=1-P$. The decomposition $f=J\xi+h$, with
$\xi=J^*f$ and $h=Qf$, is orthogonal. In particular,
\begin{equation}\label{eq:cluster-orthogonal}
 \|f\|_2^2=\|\xi\|_2^2+\|h\|_2^2.
\end{equation}
The normalization by $48$ makes $J$ an isometry. Kernels on the chamber
lattice are taken with respect to $\dd x$, so the kernel of the identity
is $48a^{-3}\mathbf1[x=x']1_4$.

Let $A$ be a connection and write
$u_A(c,c')=\exp(ie_a aA(c,c'))$ on a photon bond. The charge $e_a$ in
this notation is the charge at the level under consideration. The
phases on the internal bonds are one. Put
\[
 R_i(A)f(x)=U_A(x,x^{(i)})f(x^{(i)}),\qquad
 \mathcal L_A=\sum_{i=0}^3w_i(1-R_i(A)).
\]
The massless chamber operator is
\begin{equation}\label{eq:propagator-dirac}
 D_Af(x)=-a^{-1}\sum_{i=0}^3M_i(x)
       \big(U_A(x,x^{(i)})f(x^{(i)})-f(x)\big).
\end{equation}
At level $k$, $a=\eta_k$ and $e_a=e_k$. The reference mass is
added separately as $D_A+\mu$.
For real $A$, the involutions $R_i(A)$ are unitary and self-adjoint,
and the bond expansion gives
\begin{equation}\label{eq:dirac-realpart}
 \operatorname{Re}(f,D_Af)=\frac{\rho}{2a}(f,\mathcal L_Af),
 \qquad D_A^*=\gamma_5D_A\gamma_5,
 \qquad \gamma_5=\begin{pmatrix}0&1_2\\1_2&0\end{pmatrix}.
\end{equation}
Here $\gamma_5$ is Hermitian, $\gamma_5^2=1$, and it anticommutes with
each $\gamma_\mu$. Each unoriented bond contributes
$a^3w_i|U_A(x,x^{(i)})f(x^{(i)})-f(x)|_2^2/48$ to the energy.

On one cluster let
\[
 \mathcal L_{\mathrm{int}}=\sum_{i=0}^2w_i(1-R_i),
 \qquad R_iv(g)=v(gs_i).
\]
Its kernel consists of the constants: the three reflections
$s_0,s_1,s_2$ generate $O_{\mathrm{cub}}$. Denote its least positive
eigenvalue by $\lambda_{\mathrm{int}}$ and set
\[
 g_{\mathrm h}=\frac{\rho\lambda_{\mathrm{int}}}{2},\qquad
 c_* =\frac{\min(w_3,\lambda_{\mathrm{int}})}{18},\qquad
 c_{\mathrm{lh}}=\sqrt3(1+\rho).
\]
No numerical evaluation of $\lambda_{\mathrm{int}}$ is needed.

\begin{proposition}\label{prop:cluster-energy}
For real $A$ and $\xi=J^*f$,
\begin{equation}\label{eq:cluster-energy}
 (f,\mathcal L_Af)\ge\lambda_{\mathrm{int}}\|Qf\|_2^2,
 \qquad
 a^{-2}(f,\mathcal L_Af)\ge c_*\|\nabla_A\xi\|_2^2,
\end{equation}
where
$\nabla_{A,\mu}\xi(c)=a^{-1}(u_A(c,c+ae_\mu)\xi(c+ae_\mu)-\xi(c))$.
The same inequalities hold on a finite set $E$ of photon sites, with the
energy restricted to bonds whose endpoints lie over $E$ and the photon
differences restricted to bonds in $E$. Moreover,
\begin{equation}\label{eq:offdiagonal-bound}
 \|QD_AJ\xi\|_2\le c_{\mathrm{lh}}\|\nabla_A\xi\|_2,
 \qquad
 |(\xi,J^*D_AQf)|\le
 c_{\mathrm{lh}}\|\nabla_A\xi\|_2\|f\|_2.
\end{equation}
On $\operatorname{ran}Q$, for $t>-g_{\mathrm h}/a$,
\begin{equation}\label{eq:heavy-inverse}
 \|(QD_AQ+t)^{-1}\|_{2\to2}\le(t+g_{\mathrm h}/a)^{-1}.
\end{equation}
All constants are independent of $a,R$, the character, and the real
connection.
\end{proposition}

\begin{proof}
The internal bonds over $c$ give
$a^3(f(c,\cdot),\mathcal L_{\mathrm{int}}f(c,\cdot))$ and hence at
least $\lambda_{\mathrm{int}}a^3\|Qf(c,\cdot)\|^2$. Summing over
$c\in E$ proves the first bound in \eqref{eq:cluster-energy}.

For the second bound, consider a photon bond $c'=c+av$. The eight
labels with $ge_1=v$ correspond to the eight chamber bonds crossing it.
Write $\mathcal O_v(c)$ for this octet. Since their endpoints form
$\mathcal O_{-v}(c')$, the difference of the cluster means is the sum
of the mean of the eight bond differences and two means of $Qf$:
\begin{align*}
 u_A(c,c')\xi(c')-\xi(c)
 ={}&\frac18\sum_{x\in\mathcal O_v(c)}
       \big(U_A(x,x^{(3)})f(x^{(3)})-f(x)\big)\\
 &-\frac{u_A(c,c')}8\sum_{x'\in\mathcal O_{-v}(c')}Qf(x')
   +\frac18\sum_{x\in\mathcal O_v(c)}Qf(x).
\end{align*}
The square is bounded by $3/8$ times the sum of the squared terms in
these three octets. Distinct photon bonds use distinct chamber bonds;
each octet of cluster remainders occurs at most once. Consequently,
with $\mathcal E_3$ and $\mathcal E_{\mathrm{int}}$ denoting the two
parts of the energy,
\[
 a^2\|\nabla_A\xi\|_{2,E}^2
 \le18\big(w_3^{-1}\mathcal E_3+\|Qf\|_{2,E}^2\big)
 \le18\big(w_3^{-1}\mathcal E_3+
                \lambda_{\mathrm{int}}^{-1}\mathcal E_{\mathrm{int}}\big).
\]
This proves the local bound and therefore the global one.

The internal hops vanish on $J\xi$. Thus
\[
 (D_AJ\xi)(c,g)=-3(\rho-\gamma\cdot ge_1)
                   \nabla_{A,ge_1}\xi(c).
\]
Each of the six directions occurs eight times. The norm of the matrix
in this formula is $3(1+\rho)$, and the two directed differences on a
photon bond have equal norms. It follows that
$\|D_AJ\xi\|_2^2\le3(1+\rho)^2\|\nabla_A\xi\|_2^2$.
Projection gives the first bound in \eqref{eq:offdiagonal-bound}; the
second follows from \eqref{eq:dirac-realpart}, since $J$ and $Q$ commute
with $\gamma_5$.

For $Qf=f$, \eqref{eq:dirac-realpart} gives
$\operatorname{Re}(f,(QD_AQ+t)f)\ge(t+g_{\mathrm h}/a)\|f\|_2^2$.
The same lower bound holds for the adjoint. The operator has closed
range and its adjoint has zero kernel, so it is onto, also on the
infinite lattice. This proves \eqref{eq:heavy-inverse}.
\end{proof}

\subsection{The free symbol and its Schur complement}
\label{subsec:schur-symbol}

Use lattice units for the moment. Let $\mathcal D_m(p)$ be the symbol
of $D_0+m$, $0\le m\le1$, on
$V=\ell^2(O_{\mathrm{cub}};\C^4)$ with measure $48^{-1}\sum_g$:
\[
 (\mathcal D_m(p)v)(g)=m v(g)
 +\sum_{i=0}^2M_i(g)(v(g)-v(gs_i))
 +M_3(g)(v(g)-\exp(ip\cdot ge_1)v(gr_{e_1})).
\]
Let $D_{\mathrm{ll}},D_{\mathrm{lh}},D_{\mathrm{hl}},D_{\mathrm{hh}}$
be its blocks in $V=\operatorname{ran}J\oplus\operatorname{ran}Q$.
The mass is included in the diagonal blocks. Set
\[
 H_m(p)=D_{\mathrm{hh}}(p)^{-1},\qquad
 B_m(p)=D_{\mathrm{ll}}(p)-D_{\mathrm{lh}}(p)H_m(p)D_{\mathrm{hl}}(p),
 \qquad t(p)=\sum_\mu(1-\cos p_\mu).
\]
The compression of the full operator is the Wilson operator:
\[
 D_{\mathrm{ll}}(p)=m+i\gamma\cdot\sin p+\rho t(p).
\]
Its Schur complement also contains the contribution from the heavy
part. In particular, adding a mass to the chamber operator changes
$H_m$ as well as $D_{\mathrm{ll}}$.

\begin{proposition}\label{prop:schur-symbol}
There are positive constants $\theta_0,c_{\mathrm e},\kappa_{\mathrm{strip}}$, depending
on $\rho$ only, such that $H_m$ and $B_m$ are holomorphic in a
neighbourhood of $\{|\operatorname{Im}p_\mu|\le\theta_0\}$ and
$2\pi\Z^3$-periodic. On the real cube $[-\pi,\pi]^3$,
\begin{align}
 \|H_m(p)\|&\le(g_{\mathrm h}+m)^{-1},\label{eq:symbol-heavy}\\
 \|D_{\mathrm{lh}}(p)\|+\|D_{\mathrm{hl}}(p)\|
   &\le2c_{\mathrm{lh}}(2t(p))^{1/2},\label{eq:symbol-cross}\\
 B_m(p)&=m+i\gamma\cdot\sin p+E_m(p),\qquad
 \|E_m(p)\|\le K_\rho t(p),\label{eq:schur-expansion}\\
 \operatorname{Re}(v,B_m(p)v)&\ge(m+\rho c_*t(p))|v|_2^2,
                 \label{eq:schur-realpart}\\
 |B_m(p)v|_2&\ge c_{\mathrm e}(|p|_2^2+m^2)^{1/2}|v|_2,
                 \label{eq:schur-ellipticity}
\end{align}
where $K_\rho=\rho+2c_{\mathrm{lh}}^2/g_{\mathrm h}$. For every
multi-index $\nu$,
\begin{equation}\label{eq:schur-inverse-derivatives}
 \|\partial_p^\nu B_m(p)^{-1}\|
       \le C_\nu(\rho)(|p|_2^2+m^2)^{-(1+|\nu|)/2}.
\end{equation}
If $m>0$, the inverse exists also on
$|\operatorname{Im}p_\mu|\le\kappa_{\mathrm{strip}}m$, where the same bound holds
with $\operatorname{Re}p$ on the right, for
$\operatorname{Re}p\in[-\pi,\pi]^3$. At $m=0$ the real derivative
bound is asserted away from $p=0$.
\end{proposition}

\begin{proof}
Only the fourth hop depends on $p$. For $p=p'+ip''$,
\[
 \|\mathcal D_m(p)-\mathcal D_m(p')\|
       \le3(1+\rho)(\exp(|p''|_\infty)-1).
\]
Choose $\theta_0$ so that the right side is $g_{\mathrm h}/2$ when
$|p''|_\infty=\theta_0$. The heavy gap then gives
$\operatorname{Re}D_{\mathrm{hh}}(p)\ge m+g_{\mathrm h}/2$ on this
strip. Invertibility holds on a slightly larger strip, which proves
holomorphy up to the closed strip. For real $p$, the off-diagonal bounds
in \cref{prop:cluster-energy} give
$\|D_{\mathrm{lh}}\|,\|D_{\mathrm{hl}}\|
\le c_{\mathrm{lh}}(2t(p))^{1/2}$. This proves
\eqref{eq:schur-expansion}.

To retain the real-part bound after eliminating the heavy part, put
$u=Jv-H_mD_{\mathrm{hl}}v$. Then
$J^*u=v$, $Q\mathcal D_m u=0$, and
$\mathcal D_m u=JB_mv$. Therefore
\[
 \operatorname{Re}(v,B_mv)
   =\operatorname{Re}(u,\mathcal D_m u)
   =m\|u\|^2+\frac\rho2(u,\mathcal L_0(p)u)
   \ge(m+\rho c_*t(p))|v|_2^2.
\]
This proves \eqref{eq:schur-realpart}.

The ellipticity estimate uses the Dirac part near zero and the real
part away from zero. Put $t_0=(4K_\rho^2)^{-1}$. One has
$K_\rho\ge48/\lambda_{\mathrm{int}}>1$, since
$\lambda_{\mathrm{int}}\le12+24\sqrt2<48$.
Hence $t_0\le1/4$. If $t(p)\le t_0$, then
$|\sin p|_2^2\ge t(p)$ and
$\|E_m(p)\|\le\tfrac12t(p)^{1/2}$. Since
\[
 |(m+i\gamma\cdot\sin p)v|_2^2=(m^2+|\sin p|_2^2)|v|_2^2,
\]
it follows that
$|B_m(p)v|_2\ge\tfrac12(m^2+t(p))^{1/2}|v|_2$.
On $[-\pi,\pi]^3$, $t(p)\ge2|p|_2^2/\pi^2$. This proves
\eqref{eq:schur-ellipticity} in this region. If $t(p)\ge t_0$, then
\eqref{eq:schur-realpart} gives
$|B_m(p)v|_2\ge(m+\rho c_*t_0)|v|_2$.
The inequality $(|p|_2^2+m^2)^{1/2}\le m+\pi\sqrt3$ closes the
estimate with, for example,
\[
 c_{\mathrm e}=\min\left(0.22,\frac{\rho c_*t_0}{\pi\sqrt3}\right).
\]

The strip bounds give uniform bounds on the derivatives of $B_m$ by
Cauchy's formula. Differentiating $B_mB_m^{-1}=1$ and inducting on
$|\nu|$ proves \eqref{eq:schur-inverse-derivatives}. For the complex
strip, let $K_s$ bound $B_m$ on the fixed strip. On half that strip,
$\|\partial_\mu B_m\|\le2K_s/\theta_0$. Choose
\[
 \kappa_{\mathrm{strip}}=\min\left(\frac{\theta_0}{2},
                \frac{c_{\mathrm e}\theta_0}{12K_s}\right).
\]
Integration along the segment from $p'$ to $p'+ip''$ shows that the
change in $B_m$ is at most $c_{\mathrm e}m/2$. The lower bound is thus
at least $c_{\mathrm e}(|p'|_2^2+m^2)^{1/2}/2$. Differentiating the
inverse identity once more gives the asserted complex derivative bounds.
\end{proof}

\subsection{Free kernels and coercivity}\label{subsec:free-kernels}

Solving the heavy equation gives the inverse symbol
\begin{equation}\label{eq:inverse-symbol-blocks}
 \mathcal D_m^{-1}=
 \begin{pmatrix}
 B_m^{-1}&-B_m^{-1}D_{\mathrm{lh}}H_m\\
 -H_mD_{\mathrm{hl}}B_m^{-1}&
 H_m+H_mD_{\mathrm{hl}}B_m^{-1}D_{\mathrm{lh}}H_m
 \end{pmatrix}.
\end{equation}
The first term is independent of the cluster labels after insertion by
$J$ and $J^*$. Every other term has a weaker singularity at $p=0$:
$D_{\mathrm{lh}}$ and $D_{\mathrm{hl}}$ vanish there, whereas $H_m$
and all its derivatives are uniformly bounded.

\begin{proposition}\label{prop:free-chamber-kernel}
On the unit covering lattice let
\[
 G_m(c,g;c',g')=\int_{[-\pi,\pi]^3}\frac{\dd p}{(2\pi)^3}
       \exp(ip\cdot(c-c'))[\mathcal D_m(p)^{-1}](g,g').
\]
The bracket is the kernel on $V$ relative to the label measure
$48^{-1}\sum_g$, consistently with the chamber kernel convention.
It is well-defined also at $m=0$. It decomposes as
$G_m(x,x')=G_m^{\mathrm l}(c-c')+G_m^{\mathrm h}(x,x')$, where
\begin{align}
 |G_m^{\mathrm l}(n)|_2&\le C_G(1+|n|_\infty)^{-2},\label{eq:free-light}\\
 |G_m^{\mathrm l}(n+v)-G_m^{\mathrm l}(n)|_2
       &\le C_G(1+|n|_\infty)^{-3},\quad v\in\{\pm e_\mu\},\\
 |G_m^{\mathrm h}(x,x')|_2&\le C_G(1+|c-c'|_\infty)^{-3}.
\end{align}
Consequently $|G_m(x,x')|_2\le2C_G(1+|c-c'|_\infty)^{-2}$ and
\begin{equation}\label{eq:free-bond-difference}
 |G_m(x^{(i)},x')-G_m(x,x')|_2
       \le20C_G(1+|c-c'|_\infty)^{-3}.
\end{equation}
For $m>0$ these bounds have the additional factor
$\exp(-\kappa_{\mathrm{strip}}m|c-c'|_\infty)$. The constants are uniform for
$0\le m\le1$.
\end{proposition}

\begin{proof}
The Fourier estimate used here can be proved directly. If a periodic
symbol $T$ satisfies
$|\partial^\nu T(p)|\le A|p|_2^{-q-|\nu|}$, where $q<3$ and
$|\nu|\le N>3-q$, then its Fourier coefficients are bounded by
$C A(1+|n|_\infty)^{-(3-q)}$. For $n\ne0$, put
$r=|n|_\infty^{-1}$ and split the integral using a cutoff equal to
one for $|p|_2\le r$ and zero for $|p|_2\ge2r$. The inner integral
is bounded by $CA r^{3-q}$. On the complement integrate $N$ times in
a coordinate with $|n_\mu|=|n|_\infty$. The differentiated integrand
has integral at most $CA r^{3-q-N}$; the prefactor is $r^N$.
Derivatives of the cutoff are supported in $r\le|p|_2\le2r$ and have
the same bound. Boundary terms cancel by periodicity.

Apply this estimate with $(q,N)=(1,3)$ to $B_m^{-1}$. Multiplication
by $\exp(ip\cdot v)-1$ lowers its singularity to order zero, for which
$(q,N)=(0,4)$ applies. In \eqref{eq:inverse-symbol-blocks}, the two
off-diagonal terms have order zero and the term with two off-diagonal
factors has order $-1$. The latter can also be bounded by the
order-zero estimate. The remaining term $H_m$ is holomorphic on the
fixed strip, and its Fourier coefficients decay exponentially.
The kernel normalization costs a factor $48$ for operator-valued
symbols on $V$; this is absorbed in $C_G$.

For $m>0$, shift one coordinate of the contour by
$i\kappa_{\mathrm{strip}}m\operatorname{sgn}(n_\mu)$, where
$|n_\mu|=|n|_\infty$. The complex derivative bounds of
\cref{prop:schur-symbol} give the same Fourier estimates
on the shifted contour. This supplies the exponential factor.
The fixed-strip bound for $H_m$ is stronger than the required one.

An internal hop leaves $c$ fixed, so its difference cancels
$G_m^{\mathrm l}$. An external hop replaces $c$ by $c+ge_1$; use the
light difference bound and the two heavy bounds at these points.
The inequalities $1+|n+v|_\infty\ge(1+|n|_\infty)/2$ and
$\exp(\kappa_{\mathrm{strip}}m)<\sqrt3$ give the factor $20$ in
\eqref{eq:free-bond-difference}.
\end{proof}

The Green kernel at mesh $a$, for $0\le\mu a\le1$, is
$a^{-2}G_{\mu a}(x/a,x'/a)$. Thus
\begin{equation}\label{eq:mesh-free-kernel}
 |G(x,x')|_2\le2C_G(a+|c-c'|_\infty)^{-2},\qquad
 |G(x^{(i)},x')-G(x,x')|_2
       \le20C_Ga(a+|c-c'|_\infty)^{-3}.
\end{equation}
Fourier multiplication gives both inverse identities on finitely
supported functions. The factor $48$ in the identity kernel is
consistent with the measure $\dd x$ above.

For distinct sites put
$d_X(x,\widetilde x)=\max(|c-\widetilde c|_\infty,a)$.
The chamber graph has paths from $x$ to $\widetilde x$ of at most
$39d_X(x,\widetilde x)/a$ bonds, with photon sites in the coordinate
box spanned by $c,\widetilde c$. Indeed, the generators
$s_0,s_1,s_2$ are the simple reflections of the signed permutation
group, with nine positive roots. Their Coxeter length is at most nine
\cite[Theorem~22.14]{EtingofLieNotes}, so any two labels in one
cluster are joined by at most nine internal bonds. A monotone photon
path has at most $3|c-\widetilde c|_\infty/a$ steps. Before each step
move to a label with the required external direction, and after the
last step move to the final label. This uses at most
$10\cdot3|c-\widetilde c|_\infty/a+9\le39d_X/a$ bonds.
If $d_X$ is less than half the distance to $c'$, sum
\eqref{eq:mesh-free-kernel} along such a path. In the other case use
the two size bounds. The resulting interpolation gives, for
$0\le\alpha\le1$,
\begin{equation}\label{eq:free-holder-kernel}
 |G(x,x')-G(\widetilde x,x')|_2
 \le6240C_G d_X(x,\widetilde x)^\alpha
 \left((a+|c-c'|_\infty)^{-2-\alpha}
       +(a+|\widetilde c-c'|_\infty)^{-2-\alpha}\right).
\end{equation}

The following coercivity estimate is separate from the real-part bound.
It is needed to recover mesh-independent control of the block remainder.

\begin{proposition}\label{prop:free-coercivity}
There is $c_b=c_b(\rho)>0$ such that
\begin{equation}\label{eq:free-coercivity}
 \|D_0f\|_2^2\ge c_ba^{-2}(f,\mathcal L_0f).
\end{equation}
For a cube $E$ of side $\ell_E$, let $P_E$ be the mean over
$E\times O_{\mathrm{cub}}$. Then
\begin{equation}\label{eq:chamber-poincare}
 \|f-P_Ef\|_{2,E}^2\le C_P\ell_E^2a^{-2}(f,\mathcal L_0f)_E,
 \qquad C_P=\frac1{8c_*}+\frac1{\lambda_{\mathrm{int}}}.
\end{equation}
\end{proposition}

\begin{proof}
The Schur estimate at zero mass implies
$|B_0(p)v|\ge c_{\mathrm l}t(p)^{1/2}|v|$, with
$c_{\mathrm l}=\rho c_*/(2K_\rho)$: use the Dirac part for
$t(p)\le t_0$ and the real part for $t(p)\ge t_0$.
Write $u=Jv+h$ in a fibre and $f=\mathcal D_0(p)u$. The heavy equation
and then the light equation give
\[
 h=H_0(f_{\mathrm h}-D_{\mathrm{hl}}v),\qquad
 B_0v=f_{\mathrm l}-D_{\mathrm{lh}}H_0f_{\mathrm h}.
\]
Since $t(p)\le6$, these imply
$(2t(p))^{1/2}|v|\le A_1\|f\|$ and $\|h\|\le A_2\|f\|$, where
\[
 A_1=\frac{\sqrt2}{c_{\mathrm l}}
         \left(1+\frac{6(1+\rho)}{g_{\mathrm h}}\right),\qquad
 A_2=g_{\mathrm h}^{-1}(1+c_{\mathrm{lh}}A_1).
\]
The bond decomposition also gives
\[
 (u,\mathcal L_0(p)u)\le6t(p)|v|^2+(30+24\sqrt2)\|h\|^2.
\]
For the internal bonds the light part cancels, leaving at most
$2(w_0+w_1+w_2)\|h\|^2$. In an external bond difference, split off
the photon difference of $v$ and the two heavy endpoint values, and
bound the square by three times their squared norms. The octet count
then gives $6t(p)|v|^2+3w_3\|h\|^2$. This proves the displayed
upper bound.
This yields \eqref{eq:free-coercivity} in each fibre with
$c_b=(3A_1^2+(30+24\sqrt2)A_2^2)^{-1}$.
Fourier decomposition gives the assertion for all characters on the
torus and on the covering lattice, restoring the factor $a^{-2}$.

For \eqref{eq:chamber-poincare}, first apply the scalar lattice
Poincar\'e bound
$\|\xi-P_E\xi\|_{2,E}^2\le\ell_E^2\|\nabla_0\xi\|_{2,E}^2/8$.
In one dimension this follows by writing each difference as a sum of
nearest-neighbour differences; the coefficient of a bond at position
$j$ is $j(n-j)/2\le n^2/8$. In three dimensions use the three commuting
line averages, whose successive remainders are orthogonal.
Now $f-P_Ef=J(\xi-P_E\xi)+Qf$. Its two parts are orthogonal, and the
local bounds \eqref{eq:cluster-energy} give
\eqref{eq:chamber-poincare}. The heavy term uses $\ell_E\ge a$.
\end{proof}

